\item Suponga que un objeto de masa $M$ está suspendido de la parte baja de la cuerda de masa $m$ y longitud $L$ del problema 48.
\begin{enumerate}
    \item Demuestre que el intervalo de tiempo para que un pulso transversal recorra la longitud de la cuerda es
$$ \Delta t = 2 \sqrt{\frac{L}{mg}} \left( \sqrt{M + m} - \sqrt{M} \right) $$
    \item ¿Qué pasaría si? Demuestre que la expresión en el inciso (1) se reduce al resultado del problema 48 cuando $M = 0$.
    \item Demuestre que para $m \ll M$, la expresión en el inciso (1) se reduce a
$$ \Delta t \approx \sqrt{\frac{mL}{Mg}} $$
\end{enumerate}